% !PRINTED 428
\weberpage{428}
\begin{center}
\S\ 118. \textbf{\gesperrt{Composition of the Quadratic Forms.}}
\end{center}

Since in the algebraic theory of the class equation we must rely upon the theory of the composition of the quadratic forms, we preface the following remarks about it:

Two primitive forms $\varphi_1=(a_1,b_1,c_1)$, $\varphi_2=(a_2,b_2,c_2)$ of discriminant $-m$ are called \gesperrt{concordant} (einig, concordantes) if $a_1,a_2,\tfrac{1}{2}(b_1+b_2)$ are without a common divisor; $\varphi_1$ and $\varphi_2$ are therefore certainly concordant if $a_1,a_2$ are already without a common divisor, and since, when only the classes $k_1,k_2$ of $\varphi_1,\varphi_2$ are given, one can always choose $\varphi_1,\varphi_2$ in their classes in this way, it follows that in any two classes (which may also be identical) one can always find two concordant forms. If, then, this assumption holds, it follows from
\[
\tfrac{1}{4}(b_1^2-b_2^2) = \tfrac{1}{2}(b_1-b_2)\cdot\tfrac{1}{2}(b_1+b_2) = a_1c_1-a_2c_2,
\]
that $\tfrac{1}{2}(b_1-b_2)$ is divisible by the greatest common divisor $\delta$ of $a_1,a_2$, and from this, that the two congruences
\weq{1}{b' \equiv b_1\ (\mathrm{mod}\ 2a_1), \quad b' \equiv b_2\ (\mathrm{mod}\ 2a_2)}
can always be satisfied.

For from the first of them it follows that $b'=b_1+2\lambda a_1$, and one has therefore to determine $\lambda$ from the congruence
\[
\tfrac{1}{2}(b_1-b_2) + a_1\lambda \equiv 0\ (\mathrm{mod}\ a_2)
\]
which is always solvable. If $\mu$ is the least common multiple of $a_1$ and $a_2$, so that $a_1a_2=\mu\delta$, then by (1) $b'$ may still be replaced by $b=b'+2\mu h$, where $h$ is an arbitrary integer. Now since $\tfrac{1}{4}(b_1^2+m)$ is divisible by $a_1$, and $\tfrac{1}{2}(b_2^2+m)$ by $a_2$, % [sic] printed with 1/4 and 1/2 respectively, inconsistently
$\tfrac{1}{4}(b'^2+m)$ is divisible by $\mu$, and if we determine $h$ from the congruence
\[
\frac{b'^2+m}{4\mu} + b'h \equiv 0\ (\mathrm{mod}\ \delta)
\]
which is always solvable, since $b'$ is by assumption relatively prime to $\delta$, then $b^2+m=4a_1a_2c$ becomes divisible by $4a_1a_2$. The two classes $k_1,k_2$ are then represented by
\[
\varphi_1=(a_1,b,a_2c), \quad \varphi_2=(a_2,b,a_1c),
\]
and the form
\[
\varphi = \varphi_1\varphi_2 = (a_1a_2,b,c)
\]
%CONT>>
